% Seccion 2 de la rubrica: problematica y objetivo. Fija el objetivo
% estadistico, las tecnicas y las variables. Las cifras que se citan aqui
% vienen del analisis exploratorio; se referencian con \cref, no se repiten.

\section{Problemática y objetivo}
\label{sec:objetivo-tecnicas}

\subsection{Problemática}

El AMM acumula tres décadas de ozono al alza (\cref{sec:objetivo}) y su red
de monitoreo produce, cada hora y en 15 puntos, las concentraciones de los
precursores y las condiciones meteorológicas que gobiernan su formación. La
pregunta que este trabajo responde es qué combinación de esas condiciones
acompaña a un día en que el ozono supera la norma, y con qué certeza puede
identificarse ese día a partir de ellas. Dos obstáculos la vuelven un
problema multivariado y no una comparación de umbrales: los precursores
están fuertemente correlacionados entre sí (\ce{NO}, \ce{NO2} y \ce{NOx} son
la misma emisión medida tres veces) y la respuesta es rara con el umbral
vigente, de modo que una regla trivial acierta casi siempre sin decir nada.

\subsection{Objetivo del análisis multivariado}
\label{sec:obj:objetivo}

\textbf{Objetivo principal.} Clasificar cada registro día-estación como
excedencia o no excedencia del MDA8 en el umbral de 51 ppb, a partir de los
precursores y variables meteorológicas del mismo día, e identificar qué
condiciones latentes explican esa clasificación.

\textit{Criterio de logro:} área bajo la curva ROC $\geq 0.75$ en una
partición de validación no vista durante el ajuste, con sensibilidad y
especificidad reportadas por separado. No se utiliza la exactitud: con el
umbral vigente de 65 ppb las excedencias son 11.25\,\% de los días-estación
válidos y un clasificador constante alcanzaría 88.75\,\% de exactitud sin
identificar el fenómeno (\cref{sec:eda-cualitativas}). Por esa razón el
escenario principal usa 51 ppb, la meta a 5 años de la NOM-020-SSA1-2021
(\cref{tab:nom-limites}), con 31.96\,\% de excedencias \parencite{saito2015}.

Las estaciones no se agrupan por sus niveles medios de ozono, que la
estación explica en apenas 5.3\,\% de la varianza del MDA8, sino que se
homologan en las siete zonas de \cref{sec:contexto-estaciones}, cuyo
contenido espacial sí se sostiene en los datos.

\subsection{Técnicas estadísticas}
\label{sec:obj:tecnicas}

Se combinan técnicas de interdependencia, que reducen los predictores a ejes
interpretables, con técnicas de dependencia, que relacionan esos ejes con la
excedencia:

\begin{itemize}
  \item \textbf{Análisis de componentes principales (ACP)}: resuelve la
    multicolinealidad de los predictores y produce ejes ortogonales con
    lectura física.
  \item \textbf{Análisis factorial exploratorio (AF)}: separa la varianza
    común de la específica y expone redundancias que el ACP absorbe sin
    señalar. Sus puntajes son los predictores de los modelos de
    clasificación.
  \item \textbf{Regresión logística}: modela la probabilidad de excedencia y
    el aporte de cada predictor como razón de momios, sin supuestos de
    normalidad; ya se ha aplicado a la excedencia del MDA8 a partir de
    variables meteorológicas \parencite{otero2016}.
  \item \textbf{Análisis discriminante lineal}: contrasta a la logística bajo
    los supuestos de normalidad multivariada y covarianzas iguales
    \parencite{hosmer2013}. Box-Cox dejó los sesgos en $|\text{sesgo}|<0.5$
    (\cref{tab:sesgos-transformacion}), lo que aproxima el primero.
  \item \textbf{Modelo autorregresivo logístico}: incorpora la dependencia
    temporal de la respuesta, que la caracterización como serie de tiempo
    muestra que persiste tras remover la estacionalidad.
\end{itemize}

Todos los clasificadores se evalúan con el área bajo la curva ROC, no con
exactitud, por el desbalance señalado \parencite{saito2015}. Se requiere la
estandarización de \cref{sec:transformaciones}: en unidades originales la
presión, con valores cercanos a 710, dominaría las distancias frente al
\ce{CO}, cuyos valores son menores que 10.

\subsection{Variables y sus características}
\label{sec:obj:variables}

Las definiciones, unidades y faltantes se encuentran en el diccionario de
datos (\cref{sec:diccionario}); la \cref{tab:variables-tecnica} presenta el
papel de cada grupo en el objetivo.

\begin{table}[t]
  \centering
  \caption{Papel de cada grupo de variables en el objetivo.}
  \label{tab:variables-tecnica}
  \footnotesize
  \begin{tabular}{p{3.4cm} p{4.4cm}}
    \toprule
    Grupo & Papel \\
    \midrule
    \ce{O3}, MDA8 & Definen la clase; no son predictores \\
    \addlinespace[2pt]
    \ce{NO}, \ce{NO2}, \ce{NO_x}, \ce{CO} & Predictores, ventana matutina
      06--09 h \parencite{abdulwahab2005} \\
    \addlinespace[2pt]
    \ce{PM10}, \ce{PM2.5}, \ce{SO2} & Predictores, media diaria \\
    \addlinespace[2pt]
    \texttt{TOUT}, \texttt{RH}, \texttt{SR}, \texttt{PRS}, \texttt{WSR},
      \texttt{viento\_u}, \texttt{viento\_v} & Predictores, resúmenes
      diarios (máximo, acumulado, media, mínimo) \\
    \addlinespace[2pt]
    \texttt{temporada}, \texttt{tipo\_dia}, \texttt{zona} & Predictores
      categóricos, codificados en $k-1$ indicadoras \\
    \bottomrule
  \end{tabular}
\end{table}

Los precursores se resumen en la ventana 06--09 h porque es la emisión
matutina la que alimenta la fotoquímica del mediodía; la radiación solar se
conserva para distinguir la atenuación por aerosol de una diferencia en la
disponibilidad de luz. Las categóricas entran con $k-1$ indicadoras y una
categoría de referencia explícita (Noreste, día laborable e invierno);
incluir los $k$ niveles generaría colinealidad perfecta con el término
independiente. El ACP y el AF excluyen estas indicadoras porque las columnas
binarias inflarían artificialmente la varianza explicada.
